\scanpage{579}

\mspage{488}

b) $T$ est accessible dans $M$

b$'$) $T$ est $\underline L$-stable dans $M$.

\uline{Corollaire} 4.16.7.10. Soit $M$ une catégorie accessible. Les conditions suivantes sont équivalentes

a) Il existe une sous-catégorie fermée $\deleted{\unsure{\underline S|S}}$ de $M$, non accessible dans $M$.

b) Il existe une sous-catégorie ouverte $\deleted{\unsure{\underline T|\underline S}}$ de $M$, non accessible dans $M$.

% slanted
\margin{Exemple où ces conditions sont satisfaites: $M=k\text{-algèbres}$, $k$ un corps \unsure{comm.} infini, $\unsure{\mathcal U|U}$ étant le plus petit univers tel que $\unsure{\mathbb N|N}\in\mathcal U$.}

c) Il existe une suite transfinie \unsure{$x$}

\[\underline x=(x_\alpha)_{\alpha\in I}\]

d'objets de $M$, indexée par l'ens $I$ des ordinaux de l'Univers de référence $\mathcal U$, telle que l'on ait

\[\operatorname{Hom}(x_\alpha,x_\beta)=\varnothing\qquad\text{pour }\alpha<\beta .\]

En effet, l'équivalence de a) et b) est claire par le corollaire précédent, \deleted{\unsure{la précédente|ce qui précède}} et celle de a) et c) n'est autre que 4.16.7.4.

\uline{Démonstration de} 4.16.7.8. Nous aurons à utiliser le

\uline{Lemme} 4.16.7.11. Soient $M$ une catégorie accessible, $\pi$ un cardinal utile pour $M$, $\underline S$ une sous-catégorie pleine de $M$, $\underline S_{(\pi)}\overset{\mathrm{déf}}{=}\underline S\cap M_\pi$, d'où un foncteur pleinement fidèle

\[\varphi:\mathrm{Ind}_\pi(\underline S_{(\pi)})\longrightarrow M .\]

% archive sheet number: 578
